% The progressive-resolution axis: where the reduction happens, what the four
% schedules are, and which operator does the shrinking.
% Schedules and window lists come from data.tex, generated by
% scripts/fig_schema_data.py out of scripts/campaign_manifest.py and
% continuation/campaign_ops.py.  Nothing here is transcribed by hand.
\providecommand{\resDir}{}
\input{\resDir data}

\definecolor{bcblue}{HTML}{1F77B4}
\definecolor{bcred}{HTML}{D62728}
\definecolor{bcgreen}{HTML}{2CA02C}
\definecolor{bcorange}{HTML}{FF7F0E}
\definecolor{bcgrey}{HTML}{7A7A7A}

\begin{tikzpicture}[
  font=\footnotesize,
  lab/.style={font=\scriptsize, text=black!75},
  expl/.style={font=\scriptsize, text=black!85, anchor=west},
  panel/.style={font=\bfseries\small},
  flow/.style={-{Straight Barb[length=1.5mm,width=1.4mm]}, line width=0.5pt,
               black!70},
  gate/.style={rounded corners=1pt, draw=black!55, line width=0.5pt,
               fill=black!4, minimum height=8mm},
  cut/.style={rounded corners=1pt, draw=bcred, line width=0.8pt,
              fill=bcred!8, minimum height=8mm, text=bcred},
  opnode/.style={rounded corners=1.5pt, draw=black!55, fill=black!5,
                 line width=0.5pt, minimum height=5.5mm, minimum width=6.5mm},
  cutop/.style={rounded corners=1.5pt, draw=bcred, line width=0.8pt,
                fill=bcred!8, minimum height=6.5mm, minimum width=9mm,
                text=bcred},
  connection/.style={-{Straight Barb[length=1.6mm,width=1.5mm]},
                     line width=0.6pt, draw=black!60},
]
\providecommand{\midarrow}{}

% =====================================================================
% PANEL A -- the two reduction paths, drawn isometrically (generated)
% =====================================================================
\node[panel, anchor=west] at (-4.6,4.2) {A};
\node[anchor=west, font=\small] at (-4.2,4.2)
  {Where the resolution is reduced};
\node[anchor=west, font=\scriptsize, text=black!85] at (-4.2,3.5)
  {Each arm uses exactly one gate.  Both are real spatial changes: nothing
   is upsampled back,};
\node[anchor=west, font=\scriptsize, text=black!85] at (-4.2,3.0)
  {and $r=32$ is an exact bypass.  Drawn at $r=16$, the low leg of
   \texttt{Rprog}.};
\input{\resDir panelA}

\begin{scope}[shift={(-4.6,-6.7)}]
\begin{scope}[x=1mm, y=1mm]
% =====================================================================
% PANEL B -- the four schedules, and the sigma schedule beneath them
% =====================================================================
\def\XZ{30}      % x of epoch 0
\def\SC{4.0}     % mm per epoch
\newcommand{\ex}[1]{\XZ+\SC*(#1)}

\node[panel, anchor=west] at (0,-32) {B};
\node[anchor=west] at (6,-32) {The four resolution schedules, over \nEpochs{} epochs};
\node[expl] at (6,-37)
  {Darker means a smaller map.  Every arm is back at $32\times32$ from epoch
   $12$, well before training ends.};

% #1 row label, #2 the segment-list macro, #3 the row's y.
% \edef flattens #2 to its literal list so the inner \foreach can parse it.
\newcommand{\resrow}[3]{%
  \node[lab, anchor=east] at (28,#3) {#1};
  \edef\seglist{#2}%
  \foreach \a/\b/\v in \seglist {
    \pgfmathsetmacro{\shd}{(32-\v)/16*46 + 7}
    \fill[bcblue!\shd, draw=black!35, line width=0.3pt]
      ({\ex{\a}},#3-2.9) rectangle ({\ex{\b}},#3+2.9);
    \node[font=\scriptsize, text=black!80] at ({\ex{(\a+\b)/2}},#3) {$\v$};
  }%
}
\resrow{\texttt{R32}}{\segsRTwo}{-46}
\resrow{\texttt{Rprog}}{\segsRProg}{-54}
\resrow{\texttt{Rgentle}}{\segsRGentle}{-62}
\resrow{\texttt{Rreverse}}{\segsRReverse}{-70}

% the sigma schedule on the same epoch axis
\node[lab, anchor=east] at (28,-82) {$\sigma$, \texttt{Gplateau}};
\foreach \a/\b/\v in \segsSigma {
  \pgfmathsetmacro{\shd}{\v*58 + 5}
  \fill[bcred!\shd, draw=black!35, line width=0.3pt]
    ({\ex{\a}},-84.9) rectangle ({\ex{\b}},-79.1);
  \ifdim \v pt>0pt
    \node[font=\tiny, text=black!80] at ({\ex{(\a+\b)/2}},-82) {$\v$};
  \else
    \node[font=\tiny, text=black!80] at ({\ex{(\a+\b)/2}},-82)
      {$\sigma=0$: exact identity};
  \fi
}

% epoch axis and the three moments that matter
\foreach \e in {0,6,12,21,30} {
  \draw[black!45, line width=0.3pt] ({\ex{\e}},-88) -- ({\ex{\e}},-90);
  \node[lab, below] at ({\ex{\e}},-90) {\e};
}
\node[lab, anchor=west] at ({\ex{30}+3},-90) {epoch};
\foreach \e in {6,12,\bypassEpoch} {
  \draw[black!40, dashed, dash pattern=on 0.8mm off 0.8mm, line width=0.3pt]
    ({\ex{\e}},-88) -- ({\ex{\e}},-43);
}

% =====================================================================
% PANEL C -- the operator that does the shrinking
% =====================================================================
\node[panel, anchor=west] at (0,-100) {C};
\node[anchor=west] at (6,-100) {The two reduction operators};
\node[expl] at (6,-105)
  {\texttt{*\_bilinear} is an antialiased resample; \texttt{*\_max} is adaptive
   max pooling, whose windows are not always disjoint.};

\node[lab, anchor=west] at (6,-113)
  {\texttt{*\_bilinear}: \texttt{F.interpolate(..., mode="bilinear",
   align\_corners=False, antialias=True)}};
\node[lab, anchor=west] at (6,-119)
  {\texttt{*\_max}: \texttt{F.adaptive\_max\_pool2d}};

\node[lab, anchor=west] at (6,-127)
  {$32\to16$: windows all of width \winSixteenWidths, \winSixteenOverlaps\
   overlapping pairs -- disjoint.};
\node[lab, anchor=west] at (6,-133)
  {$32\to24$: windows all of width \winTwentyFourWidths, but
   \winTwentyFourOverlaps\ consecutive pairs \emph{overlap}:};

% one row of input cells with the first eight 32->24 windows bracketed under it
\def\CZ{104}     % x of input index 0
\def\CW{5.0}     % mm per input cell
\node[lab, anchor=west] at (\CZ,-123)
  {input index $\to$ first eight output windows};
\foreach \i in {0,...,11} {
  \draw[draw=black!45, fill=black!5, line width=0.3pt]
    ({\CZ+\CW*\i},-131) rectangle ({\CZ+\CW*(\i+1)},-126.2);
  \node[font=\tiny, text=black!70] at ({\CZ+\CW*(\i+0.5)},-128.6) {\i};
}
\foreach \a/\b [count=\j from 0] in \winTwentyFourList {
  \pgfmathsetmacro{\yy}{-133.0 - 2.4*mod(\j,3)}
  \draw[bcred, line width=1.0pt, line cap=round]
    ({\CZ+\CW*\a+0.5},\yy) -- ({\CZ+\CW*\b-0.5},\yy);
}
\node[lab, text=bcred, anchor=west] at (\CZ,-142)
  {consecutive windows share an input column};

\end{scope}
\end{scope}
\end{tikzpicture}
